\section{Fixed-base linear depth for Pisot numeration}
\label{sec:pisot-linear-depth}

We begin with the fixed-system theorem.  A Pisot linear numeration system is
a sequence $U=(u_j)_{j\ge0}$ of positive integers
which satisfies the integral recurrence whose characteristic polynomial is
the minimal polynomial of a Pisot number; the initial values need not be the
standard ones.  Here ``standard'' has its technical meaning
$u_{-1}=\cdots=u_{-d+1}=0$, $u_0=1$; we do not use it as a synonym for the
larger class considered below.  We assume throughout this section that $U$
is strictly increasing and $u_0=1$.  Let
\[
D_U=\{0,\ldots,d_U\}
\]
be its canonical greedy digit set.  The canonical greedy $U$-normal words of
length $m$ are
\begin{equation}
X_{U,m}=\left\{x\in D_U^m:
 \sum_{j=0}^k x_ju_j<u_{k+1}\quad(0\le k<m)\right\}.
\label{eq:pisot-canonical-words}
\end{equation}
Thus
\[
\Val_{U,m}(x)=\sum_{j=0}^{m-1}x_ju_j
\]
maps $X_{U,m}$ bijectively onto $\{0,\ldots,u_m-1\}$.  Define
\begin{equation}
\Fold_{U,m}(w)=
(\Val_{U,m}|_{X_{U,m}})^{-1}
       (\Val_{U,m}(w)\bmod u_m)
\label{eq:pisot-canonical-fold}
\end{equation}
and let $\Phi_{U,m}$ be the corresponding sliding code on
$D_U^{\mathbb Z}$.  Define $\ell_{\rm cau}(U,m)$ to be the least $L\ge1$ for
which the raw digit at time zero is determined, on the image of
$\Phi_{U,m}$, by the output windows at times $0,\ldots,L-1$; set it to
$\infty$ if no such $L$ exists.  Thus $\ell_{\rm cau}(U,m)-1$ is the least
anticipation of an inverse constrained to have memory zero.

The scope is not a disguised standard-initial-value convention.  For
example, standard $k$-bonacci initial conditions begin $1,1,\ldots$ and fail
strict increase, whereas the customary shift beginning $1,2,\ldots$ belongs
to the present class.  Nor do we assume Condition~F: that would be a genuine
restriction, since the principal root of $x^3-3x^2+2x-1$ is Pisot but does
not satisfy Condition~F \cite[Theorem~3]{Akiyama2000}.  The equivalence of
Condition~F with zero preservation is proved in
\cite{CartonSudberyYassawi2026}, but neither property is used below.

Put $\Delta_U=[-d_U,d_U]\cap\mathbb Z$ and, for $L\ge1$, set
\begin{equation}
\mathcal N_{U,m,L}=\left\{e\in\Delta_U^{m+L-1}:
 e_0\ne0,\quad
 \sum_{j=0}^{m-1}u_je_{t+j}\equiv0\pmod {u_m}
 \ (0\le t<L)\right\}.
\label{eq:pisot-bad-sliding-set}
\end{equation}
Every row in this system has a uniquely defined carry $k_t\in\mathbb Z$:
\begin{equation}
\sum_{j=0}^{m-1}u_je_{t+j}=k_tu_m.
\label{eq:pisot-window-carry}
\end{equation}
Equivalently,
\begin{equation}
[e_t,e_{t+1},\ldots,e_{t+m-1},-k_t]_U=0,
\label{eq:pisot-zero-word-bridge}
\end{equation}
where the last coordinate has weight $u_m$.  This is the promised bridge to
the standard language of bounded zero representations; it is not an
identification with numerical $\beta$-normalization.

\begin{lemma}[Exact collision quotient for $U$]
\label{lem:pisot-exact-collision-quotient}
For every $m\ge2$,
\begin{equation}
\ell_{\rm cau}(U,m)
=\min\{L\ge1:\mathcal N_{U,m,L}=\varnothing\}.
\label{eq:pisot-exact-depth}
\end{equation}
Let $\Gamma_{U,m}$ have vertex set $\Delta_U^{m-1}$ and an edge
$(v_0,\ldots,v_{m-2})\to(v_1,\ldots,v_{m-2},a)$ exactly when
\[
\sum_{j=0}^{m-2}u_jv_j+u_{m-1}a\equiv0\pmod {u_m}.
\]
If $B_{U,m}$ is the set of vertices with first coordinate nonzero, then
\begin{equation}
\ell_{\rm cau}(U,m)=1+\sup\{\text{edge length of a path from }B_{U,m}\}.
\label{eq:pisot-longest-path}
\end{equation}
Moreover, $\Phi_{U,m}$ is injective if and only if no directed cycle of
$\Gamma_{U,m}$ is reachable from $B_{U,m}$.
\end{lemma}

\begin{proof}
Two folded windows agree exactly when the difference of their raw values is
divisible by $u_m$, because the canonical normal words represent
$0,\ldots,u_m-1$ bijectively.  Every $e\in\Delta_U$ is realized as a
difference of two digits by taking its positive and negative parts in
$D_U$.  Hence two output $L$-blocks with distinct initial raw digits are
equivalent to a vector in $\mathcal N_{U,m,L}$, with no lost collisions.
Reading consecutive difference windows gives precisely the paths from
$B_{U,m}$ in $\Gamma_{U,m}$, proving \eqref{eq:pisot-exact-depth} and
\eqref{eq:pisot-longest-path}.  A reachable nonzero cycle can be repeated
periodically and is realized by two raw points.

It remains only to note why injectivity rules out the zero loop as a
reachable cycle.  If $d_U\ge u_m$, the constant nonzero difference sequence
$e_t=u_m$ belongs to $\Delta_U^{\mathbb Z}$ and every window sum is divisible
by $u_m$; hence $\Phi_{U,m}$ is noninjective and the difference graph has a
nonzero loop.  In the only remaining case, $d_U<u_m$.  If a vertex has an
edge to the zero vertex, its overlap coordinates vanish and its first
coordinate $a$ satisfies $a\equiv0\pmod {u_m}$.  The bound $|a|\le d_U<u_m$
then gives $a=0$, so the zero vertex is its own unique predecessor.  This
proves the last assertion without imposing a restriction on the allowed
nonstandard initial values.
\end{proof}

We next make the contraction constant explicit.  Let
$\beta=\beta_0,\beta_1,\ldots,\beta_{s-1}$ be the roots of the minimal
polynomial, with $\beta>1$ and $|\beta_i|<1$ for $i\ge1$.  The recurrence and
the initial values effectively give algebraic numbers $c_i$ such that
\begin{equation}
u_n=\sum_{i=0}^{s-1}c_i\beta_i^n,
\qquad c_0>0.
\label{eq:pisot-binet}
\end{equation}
Effective root isolation gives positive algebraic bounds $b_U,A_U$ with
\begin{equation}
b_U\beta^n\le u_n\le A_U\beta^n\qquad(n\ge0).
\label{eq:pisot-effective-growth}
\end{equation}
Indeed, isolate a point after which the total conjugate contribution in
\eqref{eq:pisot-binet} is at most $c_0\beta^n/2$, and inspect the remaining
finite prefix.

Set
\begin{equation}
K_U^{*}=
\left\lceil\frac{d_UA_U}{b_U(\beta-1)}\right\rceil.
\label{eq:pisot-analytic-carry-bound}
\end{equation}
Equation \eqref{eq:pisot-effective-growth} shows directly, without a
normalization automaton or a zero-preservation hypothesis, that every carry
in \eqref{eq:pisot-window-carry} lies in the integer interval
$[-K_U^{*},K_U^{*}]$, since
\[
|k_t|\le d_U\frac{\sum_{j<m}u_j}{u_m}
<\frac{d_UA_U}{b_U(\beta-1)}.
\]

For nonintegral $\beta$, define the effective positive separation
\begin{equation}
\delta_U^{*}=\min\left\{|a+k-\beta l|:
 a\in\Delta_U,\ |k|,|l|\le K_U^{*},\ a+k-\beta l\ne0\right\}
\label{eq:pisot-effective-separation}
\end{equation}
and
\begin{equation}
M_U=\frac1{c_0}\sum_{i=1}^{s-1}|c_i|
 \left(\frac{d_U}{1-|\beta_i|}+K_U^{*}\right),
\qquad
B_U=d_U+(1+\beta)M_U.
\label{eq:pisot-conjugate-bound}
\end{equation}
The minimum in \eqref{eq:pisot-effective-separation} is positive: if
$l\ne0$, irrationality of $\beta$ prevents equality with the integer
$a+k$, while for $l=0$ every nonzero value has modulus at least one.  All
comparisons are comparisons of algebraic numbers and are effective.
Finally let
\begin{equation}
m_U=\min\{m\ge2:\beta^m\delta_U^{*}>B_U\}.
\label{eq:pisot-contraction-threshold}
\end{equation}

For a fixed $U$, the classical finite automaton recognizing bounded zero
representations over a fixed alphabet may be used to discard interval values
that never occur \cite{Frougny1992,CartonSudberyYassawi2026}.  This can sharpen
the numerical value of $C_U$ below, but it is neither a hypothesis nor a
logical input to the proof.

\begin{lemma}[Adjacent-window Pisot contraction]
\label{lem:pisot-adjacent-window-contraction}
Suppose that $\beta$ is nonintegral, $m\ge m_U$, and two consecutive rows
of \eqref{eq:pisot-window-carry} hold.  Then
\begin{equation}
k_{t+1}=0,\qquad e_{t+m}=-k_t.
\label{eq:pisot-carry-collapse}
\end{equation}
\end{lemma}

\begin{proof}
For each embedding define
\[
A_{t,i}=\sum_{j=0}^{m-1}e_{t+j}\beta_i^j-k_t\beta_i^m.
\]
Substitution of \eqref{eq:pisot-binet} into
\eqref{eq:pisot-window-carry} gives
$\sum_i c_iA_{t,i}=0$.  For $i\ge1$,
\[
 |A_{t,i}|\le \frac{d_U}{1-|\beta_i|}+K_U^{*},
\]
and therefore $|A_{t,0}|\le M_U$.  Direct subtraction of two adjacent
dominant-embedding sums yields
\begin{equation}
\beta^m(e_{t+m}+k_t-\beta k_{t+1})
=\beta A_{t+1,0}-A_{t,0}+e_t.
\label{eq:pisot-adjacent-binet-identity}
\end{equation}
The right side has modulus at most $B_U$.  By
\eqref{eq:pisot-effective-separation} and
\eqref{eq:pisot-contraction-threshold}, the parenthesis on the left must
vanish.  Since $e_{t+m}+k_t$ and $k_{t+1}$ are integers and $\beta$ is
irrational, this is possible only in \eqref{eq:pisot-carry-collapse}.
\end{proof}

In the degree-one case set $C_U=1$.  Otherwise, for $2\le m<m_U$, let
$p_U(m)$ be the maximum edge length of a path from $B_{U,m}$ when no directed
cycle is reachable from $B_{U,m}$, and put $p_U(m)=0$ otherwise.  This finite
list is computed by exact acyclic longest-path searches.  Define
\begin{equation}
C_U=\max\left(2,
 \max_{2\le m<m_U}\left\lceil\frac{1+p_U(m)}m\right\rceil\right).
\label{eq:effective-linear-constant}
\end{equation}
Here and below the inner maximum is zero if $m_U=2$.
This is a finite, effective recipe from the recurrence and the canonical
digit set.  The large-aperture part
of the recipe is the Pisot contraction estimate, not a count of the
$(2d_U+1)^{m-1}$ difference vertices.

\begin{theorem}[Fixed-base linear causal depth and sharpness for Pisot numeration]
\label{thm:pisot-linear-depth-sharpness}
Let $U$ be any strictly increasing Pisot numeration system with $u_0=1$,
canonical greedy digit set and canonical normal words.  Then there is an
effectively computable $C_U<\infty$ such that, for every $m\ge2$,
\begin{equation*}
\Phi_{U,m}\text{ injective}
\quad\Longrightarrow\quad
\ell_{\rm cau}(U,m)\le C_Um.
\tag{A1}\label{eq:pisot-linear-A1}
\end{equation*}
Equivalently, if no first-coordinate-nonzero state of $\Gamma_{U,m}$ can
reach a directed cycle, every directed path from that initial state set has
length at most $C_Um-1$.

Let $\theta$ be the real root of
\[
x^3-2x^2+x-1
\]
and let $U_\theta=(Q_j(\theta))_{j\ge0}$ be given by
\begin{equation}
Q_0(\theta)=1,\quad Q_1(\theta)=2,\quad Q_2(\theta)=4,
\quad Q_{j+3}(\theta)=2Q_{j+2}(\theta)-Q_{j+1}(\theta)+Q_j(\theta).
\label{eq:theta-count-recurrence}
\end{equation}
For every $m\ge4$,
\begin{equation*}
\Phi_{\theta,m}\text{ is injective},\qquad
\ell_{\rm cau}(\theta,m)=2\left\lfloor\frac m2\right\rfloor-1.
\tag{A2}\label{eq:pisot-linear-A2}
\end{equation*}
More precisely, if
$\lambda_m=2\lfloor m/2\rfloor-1$, then
\begin{equation*}
\mathcal N_{\theta,m,\lambda_m-1}=\{E_m,-E_m\}\ne\varnothing,
\qquad
\mathcal N_{\theta,m,\lambda_m}=\varnothing,
\tag{A3}\label{eq:pisot-linear-A3}
\end{equation*}
where, writing concatenation by juxtaposition and changing every sign by a
prefixed minus,
\begin{equation}
\begin{gathered}
E_4=(1,-1,-1,-1,1),\\
E_m=(1,0,-E_{m-2},0,0)\quad(m\ge6\text{ even}),\\
E_m=(E_{m-1},0)\quad(m\ge5\text{ odd}).
\end{gathered}
\label{eq:theta-terminal-word}
\end{equation}
Consequently
\[
\lim_{m\to\infty}\frac{\ell_{\rm cau}(\theta,m)}m=1,
\]
so no theorem valid for every fixed system in the class of \textup{(A1)} can
replace $O_U(m)$ by $o(m)$.  This does not assert that the computed $C_U$ is
numerically optimal or that coefficient one is universal across systems.
\end{theorem}

\begin{proof}
We first prove \eqref{eq:pisot-linear-A1}.  The degree-one case is
$u_j=b^j$, $D_U=\{0,\ldots,b-1\}$; its fold is the identity and
$\ell_{\rm cau}(U,m)=1$.  Suppose that $\beta$ is nonintegral and
$m\ge m_U$.  If a vector in $\mathcal N_{U,m,L}$ has $L\ge2$, apply
Lemma~\ref{lem:pisot-adjacent-window-contraction} for
$0\le t\le L-2$.  It gives
\[
k_1=\cdots=k_{L-1}=0,\qquad
e_m=-k_0,\qquad e_{m+1}=\cdots=e_{m+L-2}=0.
\]
If $L=m+1$, the terminal difference vertex is therefore zero.  It cannot
be reached from $B_{U,m}$ when no directed cycle is reachable, because the
zero vertex has a loop.  Hence every such path has at most $m$ edges and
$\ell_{\rm cau}(U,m)\le m+1\le2m$.  The finitely many smaller apertures are
covered by the definition of $C_U$.  Lemma
\ref{lem:pisot-exact-collision-quotient} gives both formulations of
\eqref{eq:pisot-linear-A1}.

We turn to $\theta$.  The polynomial is negative at $7/4$, positive at
$351/200$, and has only one real zero; the product of the two nonreal roots
has modulus $1/\theta$.  Thus $\theta$ is Pisot.  Its defining equation
gives
\[
1=\theta^{-1}+\theta^{-2}+\theta^{-4},
\]
and every proper suffix of $1101\,0^\infty$ is lexicographically smaller.
Parry's criterion therefore gives $d_\theta(1)=1101\,0^\infty$, and the
canonical language counts are exactly \eqref{eq:theta-count-recurrence}.
In particular, the digit and difference alphabets are $\{0,1\}$ and
$\{-1,0,1\}$.  The recurrence also implies the positive recurrence
$Q_{n+4}=Q_{n+3}+Q_{n+2}+Q_n$.  Thus
$2Q_{n+1}-3Q_n>0$ follows from its four positive initial values.  Starting
with $Q_0<2Q_1$ and using
$\sum_{j<m+1}Q_j<3Q_m<2Q_{m+1}$ then gives
\begin{equation}
\sum_{j<m}Q_j(\theta)<2Q_m(\theta),
\label{eq:theta-carry-one}
\end{equation}
so every window carry belongs to $\{-1,0,1\}$.

For completeness, we verify the contraction threshold rather than invoking
an unspecified asymptotic.  If $\rho,\overline\rho$ are the nonreal roots,
write
\[
Q_n=c_\theta\theta^n+c_\rho\rho^n+\overline{c_\rho}\,\overline\rho^n.
\]
The rational isolating interval above and the residue formula for the
generating function
\[
\sum_{n\ge0}Q_nz^n=\frac{1+z^2}{1-2z+z^2-z^3}
\]
give
\[
1.26<c_\theta<1.27,\qquad |c_\rho|<0.19,\qquad |\rho|<0.756.
\]
Here $c_\alpha=(\alpha^2+1)/(3\alpha^2-4\alpha+1)$ for every root
$\alpha$, and $|\rho|=\theta^{-1/2}$.  Thus all displayed decimals are
terminating rational bounds obtained by substitution of the endpoints
$7/4$ and $351/200$ and exact rational arithmetic.
Hence the dominant-embedding error in the proof of Lemma
\ref{lem:pisot-adjacent-window-contraction} is less than $1.54$.  Here
$\delta_\theta=2-\theta>0.245$, and
\[
\theta^6\delta_\theta>7.03
>1+(1+\theta)1.54.
\]
It follows that \eqref{eq:pisot-carry-collapse} holds for every $m\ge6$.
The apertures $m=4,5$ will be included in the exact terminal check below.

It remains to classify the longest paths.  Fix $m\ge6$ and a vector with
at least two collision rows.  After the carry collapse, all its coordinates
beyond $m$ vanish except possibly $e_m=-k_0$.  Define the full suffix sums
\[
R_t=\sum_{j\ge0}Q_j(\theta)e_{t+j}.
\]
Every collision row with $t\ge1$ is now $R_t=0$, and the first row together
with $e_m=-k_0$ gives $R_0=0$.  The recurrence implies the exact identity
\begin{equation}
R_t=e_t+e_{t+2}+2R_{t+1}-R_{t+2}+R_{t+3}.
\label{eq:theta-suffix-renormalization}
\end{equation}
Thus four consecutive zero suffix sums force $e_t=-e_{t+2}$.

Let first $m$ be even and take $L=m-2=\lambda_m-1$.  We have
$R_t=0$ for $0\le t\le m-3$.  The terminal equation is
\[
e_{m-3}+2e_{m-2}+4e_{m-1}+7e_m=0.
\]
In $\{-1,0,1\}^4$ its complete solution set is
\begin{equation}
(0,0,0,0),\qquad \pm(1,1,1,-1).
\label{eq:theta-four-tail-table}
\end{equation}
The zero tail propagates backwards to the zero word.  From a nonzero tail,
the equations $R_{m-4}=R_{m-5}=0$ give respectively
$e_{m-4}=-e_{m-3}$ and $e_{m-5}=0$; then
\eqref{eq:theta-suffix-renormalization} propagates the alternating pair
uniquely to the left.  This is precisely the recursion
\eqref{eq:theta-terminal-word}, and both signs occur.  Hence the first set
in \eqref{eq:pisot-linear-A3} is $\{\pm E_m\}$.  With one more row, the
terminal equation is
\[
e_{m-2}+2e_{m-1}+4e_m=0,
\]
whose only solution in $\{-1,0,1\}^3$ is zero.  Backward propagation then
contradicts $e_0\ne0$.  Thus $\mathcal N_{\theta,m,\lambda_m}$ is empty.

Now let $m$ be odd and take $L=m-3=\lambda_m-1$.  The terminal five-digit
equation has, up to sign, exactly the following solutions:
\begin{equation}
(0,0,0,0,0),\quad
(1,1,1,-1,0),\quad
(1,0,1,1,-1),\quad
(1,-1,-1,-1,1).
\label{eq:theta-five-tail-table}
\end{equation}
Successive solution of $R_{m-5}=0$ gives respectively
$0,-1,0,0$ for the preceding digit.  Since $m-5$ is even,
\eqref{eq:theta-suffix-renormalization} propagates this digit to $e_0$ up
to sign.  The zero tail and the last two nonzero tails therefore give
$e_0=0$; only $\pm(1,1,1,-1,0)$ survive.  They give exactly the even
terminal word for aperture $m-1$ followed by zero, namely $\pm E_m$.
With one more row the terminal four-digit alternatives are
\eqref{eq:theta-four-tail-table}; now $m-5$ has even parity and its forced
zero propagates to $e_0$.  Hence the next obstruction set is empty.

For $m=4,5$, successive exact division in the two or three window
congruences gives
\[
\begin{array}{c|c|c|c}
m&Q_m&\mathcal N_{\theta,m,\lambda_m-1}
&\mathcal N_{\theta,m,\lambda_m}\\ \hline
4&12&\{\pm(1,-1,-1,-1,1)\}&\varnothing\\
5&21&\{\pm(1,-1,-1,-1,1,0)\}&\varnothing.
\end{array}
\]
This is an exhaustive integer calculation: in each row the quotient is in
$\{-1,0,1\}$ by \eqref{eq:theta-carry-one}, and the last digit determines
the preceding bounded digit successively.  Thus \eqref{eq:pisot-linear-A3}
holds for every $m\ge4$.  Lemma
\ref{lem:pisot-exact-collision-quotient} now gives
\eqref{eq:pisot-linear-A2}, including injectivity, and the limiting statement
is immediate.
\end{proof}

\begin{corollary}[Global bounded-zero strip dichotomy]
\label{cor:pisot-zero-strip-synchronization}
For a fixed strictly increasing Pisot numeration $U$ as in the theorem and an
aperture $m\ge2$,
consider all bounded zero words of length $m+1$ of the form
\[
[e_t,\ldots,e_{t+m-1},-k_t]_U=0
\]
which overlap one digit at a time.  Exactly one of the following global
alternatives holds: some strip with a nonzero difference has a finite prefix
followed by a periodic bounded zero representation; or every nonzero strip
has fewer than $(C_U+1)m$ rows.  The latter bound has sharp order $\Theta(m)$
for $U_\theta$.
\end{corollary}

\begin{proof}
The overlapping rows are exactly paths in $\Gamma_{U,m}$ by
\eqref{eq:pisot-zero-word-bridge}.  A cycle reachable from $B_{U,m}$ gives
the first alternative by repeating the cycle.  If no such cycle is
reachable, \eqref{eq:pisot-linear-A1} bounds every path from $B_{U,m}$.
Unless a strip is zero, at most $m-1$ initial shifts pass before one of its
nonzero coordinates occupies the first position, after which fewer than
$C_Um$ further rows are possible.  The two alternatives are mutually
exclusive because the periodic tail makes strips arbitrarily long.  The
words $E_m$ in \eqref{eq:pisot-linear-A3} give the matching lower order for
$U_\theta$.
\end{proof}
